\documentclass[11pt,a4paper]{article}
\usepackage[T1]{fontenc}
\usepackage[utf8]{inputenc}
\usepackage[sc]{mathpazo}
\usepackage{courier}
\usepackage[margin=19mm,top=18mm,bottom=19mm,headheight=14pt,headsep=7mm,footskip=10mm]{geometry}
\usepackage{microtype,booktabs,tabularx,array,amsmath}
\usepackage{caption,titlesec,fancyhdr}
\usepackage[hidelinks]{hyperref}
\urlstyle{same}
\captionsetup{font=small,labelfont=bf,justification=raggedright,singlelinecheck=false,skip=5pt,hypcap=false}
\titleformat{\section}{\large\bfseries}{\thesection}{0.6em}{}
\titlespacing*{\section}{0pt}{10pt}{5pt}
\setlength{\parindent}{0pt}
\setlength{\parskip}{4pt}
\setlength{\tabcolsep}{6pt}
\renewcommand{\arraystretch}{1.08}
\newcolumntype{L}{>{\raggedright\arraybackslash}X}
\newenvironment{brieftable}{\par\noindent\begin{minipage}{\linewidth}}{\end{minipage}}
\pagestyle{fancy}
\fancyhf{}
\fancyhead[L]{\footnotesize Research results brief}
\fancyhead[R]{\footnotesize Controlled RAG evaluation}
\fancyfoot[C]{\small\thepage}
\renewcommand{\headrulewidth}{0.3pt}
\raggedbottom
\hypersetup{pdftitle={RAG answer evaluation under conflicting evidence},pdfsubject={Measured results for R1--R5, C1--C5 and controls}}
% The table bodies between "% BEGIN X" and "% END X" are written by scripts/build_results.py
% from outputs/results/ (judge codex:gpt-6-luna:low, variant=main, rendering=attributed unless
% stated). Deltas are computed from unrounded values, then rounded to two decimals.
\begin{document}
\fontsize{10.5}{12.9}\selectfont
{\fontsize{18}{21}\selectfont\bfseries RAG answer evaluation under conflicting evidence\par}
\vspace{3pt}
{\large ConflictScore, RAGChecker and Q-CARE\par}
\vspace{3pt}
{\small Results brief \quad$\cdot$\quad implementation 1.0.0 \quad$\cdot$\quad 19 November 2026}

\section{Experimental setup}
\textbf{Research question.} How do published evaluators score fixed RAG answers when sources
contradict each other, have different assigned reliability, or express uncertainty?
The experiment uses a fictional person, Alex Meron, and two oracle documents supplied directly
to the evaluators; it does not evaluate an answer generator or a retrieval system.
There is no independently verified birthday.

\textbf{Questions.} \textsc{Date} (\texttt{Q\_DATE}): ``What exact date was Alex Meron born?''
\textsc{Year} (\texttt{Q\_YEAR}): ``Which year was Alex Meron born?''
Both questions use the same evidence and the same five candidate answers.

\begin{brieftable}
\captionof{table}{Fixed candidate answers, reproduced verbatim.}
\label{tab:answers}
{\small
\begin{tabularx}{\linewidth}{@{}p{7mm}p{30mm}L@{}}
\toprule
ID & Interpretation & Exact answer text \\
\midrule
R1 & Unqualified A date & Alex Meron was born on 3 March 1990. \\
R2 & Month and year & Alex Meron was born in March 1990. \\
R3 & Shared year only & Alex Meron was born in 1990. \\
R4 & Vague uncertainty & The exact date of birth is uncertain. \\
R5 & Attributed conflict & Archive A gives 3 March 1990, while Registry B gives 17 March 1990. Both agree on March 1990, but the precise day is disputed. \\
\bottomrule
\end{tabularx}}
\end{brieftable}\par\vspace{3pt}

\textbf{Evidence.} In C1, Archive A states ``Alex Meron was born on 3 March 1990.''
Registry B states ``Alex Meron was born on 17 March 1990.''
Each primary scenario retains these conflicting days and the common month and year.
Three controls vary only one factor of C1: NC (both sources give 3 March), ORD\_BA
(document order B, A) and DUP\_A3 (Archive A repeated three times under the same source).

\begin{brieftable}
\captionof{table}{Primary scenarios. Reliability is assigned synthetically, not estimated.}
\label{tab:scenarios}
{\small
\begin{tabularx}{\linewidth}{@{}p{7mm}p{26mm}p{26mm}L@{}}
\toprule
ID & A reliability & B reliability & Source wording relative to C1 \\
\midrule
C1 & High (0.9) & High (0.9) & Both sources assert a definite date. \\
C2 & High (0.9) & Low (0.3) & Unchanged; A assigned higher reliability. \\
C3 & Low (0.3) & High (0.9) & Unchanged; B assigned higher reliability. \\
C4 & High (0.9) & High (0.9) & B: ``was probably born on 17 March 1990.'' \\
C5 & High (0.9) & High (0.9) & A: ``was probably born on 3 March 1990.'' \\
\bottomrule
\end{tabularx}}
\end{brieftable}\par\vspace{3pt}

\textbf{Visible inputs.} The default document format is
\texttt{[Source: Archive A]} followed by its text, and analogously for B, in A--B order.
All three evaluators receive the candidate answer and the document texts; RAGChecker and
Q-CARE also receive the question. Reliability metadata and numeric weights are \emph{not}
passed to the scoring interfaces. A separate run inserts \texttt{assigned reliability: high/low}
inside the source prefix (\S\,3); ``probably'' is visible source text in C4/C5.

\textbf{RAGChecker references.} RAGChecker additionally receives one fixed reference per
scenario and question, shared across R1--R5 (Table~\ref{tab:refs}). These are authored,
evidence-grounded ideal answers, not verified ground truth or independent human gold scores.

\begin{brieftable}
\captionof{table}{RAGChecker reference answers (verbatim). ORD\_BA and DUP\_A3 reuse the C1 reference.}
\label{tab:refs}
{\small
\begin{tabularx}{\linewidth}{@{}>{\raggedright\arraybackslash}p{19mm}L@{}}
\toprule
Scenario & Exact reference text \\
\midrule
\multicolumn{2}{@{}l}{\textit{Exact-date question (\textnormal{\textsc{Date}})}} \\
C1 & The retrieved records disagree about Alex Meron's exact birth date. Archive A reports 3 March 1990, whereas Registry B reports 17 March 1990. Both agree on March 1990. The available evidence does not establish which day is correct. \\
C2 & The retrieved records disagree about Alex Meron's exact birth date. Archive A, the source assigned higher reliability, reports 3 March 1990, whereas Registry B, assigned lower reliability, reports 17 March 1990. Both agree on March 1990. The evidence favours 3 March 1990 but does not conclusively establish which day is correct. \\
C3 & The retrieved records disagree about Alex Meron's exact birth date. Registry B, the source assigned higher reliability, reports 17 March 1990, whereas Archive A, assigned lower reliability, reports 3 March 1990. Both agree on March 1990. The evidence favours 17 March 1990 but does not conclusively establish which day is correct. \\
C4 & Archive A reports that Alex Meron was born on 3 March 1990, whereas Registry B states that he was probably born on 17 March 1990. Both agree on March 1990. Given their equal assigned reliability, 3 March 1990 is better supported by the supplied evidence because Archive A reports it without qualification, although a conflicting tentative report remains. \\
C5 & Registry B reports that Alex Meron was born on 17 March 1990, whereas Archive A states that he was probably born on 3 March 1990. Both agree on March 1990. Given their equal assigned reliability, 17 March 1990 is better supported by the supplied evidence because Registry B reports it without qualification, although a conflicting tentative report remains. \\
NC & Alex Meron was born on 3 March 1990. \\
\addlinespace[3pt]
\multicolumn{2}{@{}l}{\textit{Year question (\textnormal{\textsc{Year}})}} \\
All & Alex Meron was born in 1990. \\
\bottomrule
\end{tabularx}}
\end{brieftable}\par\vspace{3pt}

\textbf{Metrics and code.} ConflictScore is scored with the paper's response metrics CS-C and
CS-R (Section~3 of the paper), using the paper's claim-decomposition and evidence prompts
(Appendix~E, Tables~10--11). The official repository contains no code for these metrics, so this part is
our implementation of the published definition; its tests reproduce the paper's worked examples.
Neither paper prompt contains the question, so CS-C/CS-R are identical for \textsc{Date} and \textsc{Year}.
RAGChecker 0.1.9 with RefChecker 0.2.17 (joint checking) and Q-CARE (\texttt{parsing=paper})
run their official code unchanged.

\textbf{LLM judges.} All three evaluation methods require LLM-based judgments. We used the Codex CLI
as the primary judge, with the requested GPT-6 model. To assess the sensitivity of the results to the
choice of judge, we additionally evaluated C1 using Claude Sonnet~5.5 and Claude Opus~5.5. The same five
candidate answers were used throughout; only the evaluation judge was changed.

\section{Main numerical comparison: C1}
\textbf{Metric key.} All displayed scores lie in $[0,1]$. Arrows give each metric's native
direction ($\downarrow$ lower is better, $\uparrow$ higher is better); the metrics share no common
quality scale. A dagger ($\dagger$) marks dependence on the authored reference answer.

{\small
\begin{tabularx}{\linewidth}{@{}p{44mm}L@{}}
\toprule
Metric & Definition \\
\midrule
CS-C ($\downarrow$) & Fraction of the answer's claims with at least one supporting \emph{and} at least one contradicting document. \\
CS-R ($\downarrow$) & Mean over claims of $|D^-|/(|D^+|+|D^-|)$, the contradicting share of the relevant documents; claims with no relevant document are excluded (the paper leaves this case undefined). \\
RC \texttt{faithfulness} ($\uparrow$) & Fraction of response claims entailed by at least one retrieved chunk. \\
RC reference \texttt{f1}$^\dagger$ ($\uparrow$) & Harmonic mean of reference-based precision (response claims entailed by the reference) and recall (reference claims entailed by the response); 0 if either is 0. \\
QC \texttt{verifiableness} ($\uparrow$) & Fraction of answer atomic facts supported by at least one retrieved chunk. \\
QC \texttt{completeness} ($\uparrow$) & Fraction of query sub-questions covered by the answer's facts; measures requested information, not evidential support. \\
\bottomrule
\end{tabularx}}\par\vspace{3pt}

\begin{brieftable}
\captionof{table}{C1 scores for all five answers, separated by question. CS = ConflictScore;
RC = RAGChecker; QC = Q-CARE. Values are rounded to two decimals.}
\label{tab:baseline}
{\small
\begin{tabular*}{\linewidth}{@{\extracolsep{\fill}}lrrrrrr@{}}
\toprule
 & \multicolumn{2}{c}{CS} & \multicolumn{2}{c}{RC} & \multicolumn{2}{c}{QC} \\
\cmidrule(lr){2-3}\cmidrule(lr){4-5}\cmidrule(l){6-7}
Answer & CS-C$\downarrow$ & CS-R$\downarrow$ & Faithf.$\uparrow$ & F1$^\dagger\uparrow$ & Verif.$\uparrow$ & Compl.$\uparrow$ \\
\midrule
\multicolumn{7}{@{}l}{\textit{Exact-date question (\textnormal{\textsc{Date}})}} \\
% BEGIN BASELINE_DATE
R1 & 1.00 & 0.50 & 1.00 & 0.00 & 1.00 & 1.00 \\
R2 & 0.00 & 0.00 & 1.00 & 0.40 & 1.00 & 0.00 \\
R3 & 0.00 & 0.00 & 1.00 & 0.00 & 1.00 & 0.00 \\
R4 & 0.00 & 1.00 & 0.00 & 0.40 & 0.00 & 0.00 \\
R5 & 0.00 & 0.00 & 0.75 & 1.00 & 0.67 & 1.00 \\
% END BASELINE_DATE
\addlinespace[3pt]
\multicolumn{7}{@{}l}{\textit{Year question (\textnormal{\textsc{Year}})}} \\
% BEGIN BASELINE_YEAR
R1 & 1.00 & 0.50 & 1.00 & 0.67 & 1.00 & 1.00 \\
R2 & 0.00 & 0.00 & 1.00 & 1.00 & 1.00 & 1.00 \\
R3 & 0.00 & 0.00 & 1.00 & 1.00 & 1.00 & 1.00 \\
R4 & 0.00 & 1.00 & 0.00 & 0.00 & 0.00 & 0.00 \\
R5 & 0.00 & 0.00 & 1.00 & 0.00 & 0.67 & 1.00 \\
% END BASELINE_YEAR
\bottomrule
\end{tabular*}}
\end{brieftable}\par\vspace{3pt}

\textbf{Reading the comparison.} ConflictScore flags only R1: its single claim is supported by
Archive A and contradicted by Registry B (CS-C 1.00, CS-R 0.50), for both questions. RC faithfulness
and QC verifiableness instead give R1 1.00, against 0.75/0.67 for R5 on \textsc{Date}.
R4's single claim is contradicted by both definite sources and supported by neither, so it is not
``conflicted'' (CS-C 0) but has the maximal contradiction ratio (CS-R 1.00).
R5 scores 0/0, tied with R2 and R3: each attributed claim is supported by the source it names and
judged \textsc{Irrelevant} by the other, and ``the precise day is disputed'' is \textsc{Irrelevant}
for both documents and therefore excluded from CS-R.
Reference F1 favours R5 on \textsc{Date} (1.00 versus 0.00 for R1). For \textsc{Year}, R2/R3 gain QC
completeness (0.00 $\to$ 1.00) and RC F1 (0.40/0.00 $\to$ 1.00/1.00); R5's year F1 is 0.00 because its
attributed claims have zero precision against the year-only reference, despite recall 1.00.

\section{C2--C5 and controls: evidence experiments}
\textbf{Reliability: separate evidence visibility from reference choice.}
In C2/C3 every reference-free metric, including CS-C and CS-R, keeps its C1 value for every answer
and both questions: the evaluator-visible source texts are identical.
The only changes are in reference-dependent RAGChecker metrics, caused by the scenario-specific reference.

\begin{brieftable}
\captionof{table}{Changes in exact-date RC F1 relative to C1. Holding the reference at C1
removes every change. Year F1 changes are also zero for all answers.}
\label{tab:reliability}
{\small
\begin{tabular*}{\linewidth}{@{\extracolsep{\fill}}llrrrrr@{}}
\toprule
Scenario & Reference used & $\Delta$R1 & $\Delta$R2 & $\Delta$R3 & $\Delta$R4 & $\Delta$R5 \\
\midrule
% BEGIN RELIABILITY
C2 & Scenario-specific & +0.25 & -0.15 & 0.00 & -0.40 & -0.27 \\
C2 & C1 (held fixed) & 0.00 & 0.00 & 0.00 & 0.00 & 0.00 \\
C3 & Scenario-specific & +0.44 & -0.15 & 0.00 & -0.15 & -0.27 \\
C3 & C1 (held fixed) & 0.00 & 0.00 & 0.00 & 0.00 & 0.00 \\
% END RELIABILITY
\bottomrule
\end{tabular*}}
\end{brieftable}\par\vspace{3pt}

\textbf{Control and text rendering.} With the C1 reference held fixed, all 110 RC metric
values per C2/C3 scenario (five answers, two questions, 11 metrics) equal C1 exactly.
In separate \textsc{Date} runs that write the reliability labels into the source text, every
reference-free metric, CS-C and CS-R included, is again unchanged between C1, C2 and C3 (one R5 claim,
``Both agree on March 1990'', is labelled supported by Registry B only in C2, without effect on CS-C/CS-R).
Merely adding the labels changes R5's RC faithfulness from 0.75 to 0.50 in all three scenarios:
a rendering effect, not reliability weighting. No reliability awareness is established;
invariance to metadata the interfaces never receive is expected.

\textbf{Source uncertainty.} Table~\ref{tab:hedges} lists every non-zero C4/C5 change
in the six primary metrics of Table~\ref{tab:baseline}, for both questions,
with the default rendering and each scenario's reference.

\begin{brieftable}
\captionof{table}{Hedging effects relative to C1. Unlisted primary metric/question pairs
have zero change for every answer; 0.00 denotes an unchanged score.
F1$^\dagger$ changes can reflect the changed reference as well as its claim judgments.}
\label{tab:hedges}
{\small
\begin{tabular*}{\linewidth}{@{\extracolsep{\fill}}llrrrrr@{}}
\toprule
Question & Scenario / metric & $\Delta$R1 & $\Delta$R2 & $\Delta$R3 & $\Delta$R4 & $\Delta$R5 \\
\midrule
% BEGIN HEDGES
\textsc{Date} & C4 / CS-C & 0.00 & 0.00 & 0.00 & +1.00 & +0.25 \\
 & C4 / CS-R & 0.00 & 0.00 & 0.00 & -0.50 & +0.25 \\
 & C4 / RC F1$^\dagger$ & +0.33 & -0.07 & 0.00 & -0.40 & -0.25 \\
 & C4 / QC verifiableness & 0.00 & 0.00 & 0.00 & +1.00 & -0.33 \\
 & C5 / CS-C & 0.00 & 0.00 & 0.00 & +1.00 & 0.00 \\
 & C5 / CS-R & 0.00 & 0.00 & 0.00 & -0.50 & 0.00 \\
 & C5 / RC faithfulness & -1.00 & 0.00 & 0.00 & 0.00 & -0.50 \\
 & C5 / RC F1$^\dagger$ & +0.50 & -0.40 & 0.00 & -0.40 & -0.50 \\
 & C5 / QC verifiableness & -1.00 & 0.00 & 0.00 & +1.00 & -0.33 \\
\addlinespace[3pt]
\textsc{Year} & C4 / CS-C & 0.00 & 0.00 & 0.00 & +1.00 & +0.25 \\
 & C4 / CS-R & 0.00 & 0.00 & 0.00 & -0.50 & +0.25 \\
 & C4 / QC verifiableness & 0.00 & 0.00 & 0.00 & +1.00 & -0.33 \\
 & C5 / CS-C & 0.00 & 0.00 & 0.00 & +1.00 & 0.00 \\
 & C5 / CS-R & 0.00 & 0.00 & 0.00 & -0.50 & 0.00 \\
 & C5 / RC faithfulness & -0.50 & 0.00 & 0.00 & +1.00 & -0.25 \\
 & C5 / QC verifiableness & -1.00 & 0.00 & 0.00 & +1.00 & -0.33 \\
% END HEDGES
\bottomrule
\end{tabular*}}
\end{brieftable}\par\vspace{3pt}

A hedged source still \textsc{Supports} a definite claim: in C5, R1 keeps CS-C 1.00 and CS-R 0.50,
while its exact-date RC faithfulness and QC verifiableness fall from 1.00 to 0.00.
In both hedge scenarios the hedged source supports R4's ``uncertain'', so R4 becomes conflicted
(CS-C 1.00, CS-R 0.50) and gains QC verifiableness, but keeps QC completeness 0.00.
In C4 the judge labels R5's claim ``Archive A gives 3 March 1990'' as contradicted by the hedged
Registry B (CS-C/CS-R 0.25), a label not produced in C1 or C5. R5's exact-date RC faithfulness
is 0.75 in C4 and 0.25 in C5. Both asymmetries rest on single judge labels and are not
demonstrated systematic treatments of the sources.
The C4/C5 references name the unhedged date as better supported (Table~\ref{tab:refs}).
Against them, exact-date RC F1 ranks R5 first in C4 (0.75; R1 and R2 0.33; R3 and R4 0.00),
but in C5 R1 ties with R5 (0.50) although R1 asserts the hedged, less-supported date.
The judge extracts a plain ``born on 3 March 1990'' claim from the C4 reference, but from the C5
reference only Registry B's report of 17 March, and it judges R1 to entail Archive A's tentative
3 March report. No candidate answer states that 17 March is better supported.

\begin{brieftable}
\captionof{table}{ConflictScore CS-C\,/\,CS-R ($\downarrow$) in the controls (identical for both questions).
C1 for comparison: R1 1.00\,/\,0.50, R2 and R3 0\,/\,0, R4 0\,/\,1.00, R5 0\,/\,0.}
\label{tab:controls}
{\small
\begin{tabular*}{\linewidth}{@{\extracolsep{\fill}}lccccc@{}}
\toprule
Control & R1 & R2 & R3 & R4 & R5 \\
\midrule
% BEGIN CONTROLS
NC & 0.00 / 0.00 & 0.00 / 0.00 & 0.00 / 0.00 & 0.00 / 1.00 & 0.00 / 0.33 \\
ORD\_BA & 1.00 / 0.50 & 0.00 / 0.00 & 0.00 / 0.00 & 0.00 / 1.00 & 0.00 / 0.00 \\
DUP\_A3 & 1.00 / 0.25 & 0.00 / 0.00 & 0.00 / 0.00 & 0.00 / 1.00 & 0.00 / 0.00 \\
% END CONTROLS
\bottomrule
\end{tabular*}}
\end{brieftable}\par\vspace{3pt}

\textbf{Controls.} Reversing the document order changes no metric of any evaluator.
Without a conflict (NC), R1 is no longer conflicted, and R5's now-false statement about Registry B
is detected (CS-R 0.33); RC faithfulness and QC verifiableness for R5 also fall (0.75 $\to$ 0.50,
0.67 $\to$ 0.33 on \textsc{Date}). Repeating Archive A three times lowers R1's CS-R from 0.50 to 0.25:
ConflictScore counts duplicates as independent evidence. RAGChecker is unaffected (max over chunks),
and Q-CARE's answer-independent C-Prec@10 rises from 0.20 to 0.40.

\section{Findings and limitations}
\textbf{Contradiction and support measure different things.}
ConflictScore uses both supporting and contradicting labels and flags R1's contested day.
RC faithfulness takes the maximum entailment across chunks and QC verifiableness the union of
supported facts, so their 1.00 for R1 means that \emph{some} source supports it, with nothing
subtracted for the other source's contradiction.

\textbf{Overconfident R1 versus conflict-aware R5.}
No reference-free metric rewards R5's acknowledgement of the conflict. CS-C/CS-R are penalty-only:
R5 avoids a penalty because each attributed claim is checked only against the source it names,
and its statement of the dispute counts as irrelevant evidence. RC verifies three of R5's four
extracted claims (``the precise birth day is disputed'' is Neutral against each chunk individually);
QC verifies the two source-specific facts but not the combined agreement/disagreement fact (2/3).
Only RAGChecker's reference-based F1 favours R5, and it measures agreement with our authored
ideal answer, not independent confirmation of answer quality.

\textbf{Specificity and cautious wording.}
ConflictScore never sees the question, so R2 and R3 score 0/0 even when the exact date is asked;
QC completeness distinguishes the requested granularity. R4's uncertainty is contradicted by each
definite source individually: CS-R 1.00, and 0.00 for RC faithfulness and both QC metrics in C1.
A good native score is not a good overall answer-quality judgment.

\textbf{Cross-judge check (C1; Sonnet both questions, Opus \textsc{Date}).} R1--R4 keep their RC
faithfulness and QC verifiableness under both Claude judges, but R5 is judged less grounded: QC
verifiableness 0.67 $\to$ 0.50 under both, and RC faithfulness 0.75 $\to$ 0.50 (\textsc{Date}) and
1.00 $\to$ 0.50 (\textsc{Year}) under Sonnet. Reference-based metrics are less stable: RC precision for R1 falls from 1.00 to 0.00
under both Claude judges, and R5's \textsc{Date} QC completeness from 1.00 to 0.00.
CS-C and CS-R are identical under all three judges for every answer, although the claim wording
differs and Opus splits R2 into two claims. All Claude-judge results are scored.

\textbf{Limits of inference.} These controlled synthetic experiments use adapted implementations
and a substituted judge; they are not reproductions of the papers' experimental configurations.
Temperature was not controllable, and each label is a single judge sample; repeated identical
prompts reuse one cached judgment rather than drawing independent samples, so judge variance is
not estimated. The served Codex model is unverified. Q-CARE's API judge is outside its documented
local-model setup. Claim extraction and entailment depend on the question and on the judge.
CS-R relies on a documented convention for claims without relevant evidence (excluded from the mean);
counting them as 0 would change C4 R5 from 0.25 to 0.125 and NC R5 from 0.33 to 0.25, and no ranking. No independent human
gold scores exist, and these results do not establish general method failure.

\vspace{2pt}
{\fontsize{8.5}{10.3}\selectfont\raggedright
\textbf{Audit trail (repository-relative).}
Tables: \path{outputs/results/comparison.csv}, \path{conflictscore_published.csv},
\path{deltas_vs_C1.csv}, \path{judge_comparison_agreement.csv};
claim-level labels: \path{outputs/results/claim_level_traces.json};
judgment sources: \path{outputs/results/judgment_provenance.csv} and \path{judge_records.jsonl}.
Runs and configurations: \path{outputs/runs/*/run_meta.json}, commands in \path{scripts/run_all.sh};
tables: \path{scripts/build_results.py}. Implementation: \path{outputs/reports/implementation_summary.pdf}.\par
}
\end{document}
